\documentclass{article}
\usepackage{amsmath,amssymb,amsthm}
\usepackage{algorithm}
\usepackage{algpseudocode}
\usepackage{hyperref}
\usepackage{enumitem}
\newtheorem{theorem}{Theorem}
\newtheorem{lemma}{Lemma}
\newcommand{\E}{\mathbb{E}}
\newcommand{\R}{\mathbb{R}}
\newcommand{\norm}[1]{\left\|#1\right\|}
\newcommand{\inner}[1]{\left\langle #1 \right\rangle}
\begin{document}

\section*{Algorithm Name}
\textbf{Error‑Feedback Stochastic Gradient Descent (EF‑SGD)}  

The method augments a stochastic gradient step with a memory vector that stores the quantization error from the previous iteration and adds it to the current stochastic gradient before quantization.

\section*{Mathematical Formulation}
Let $f:\R^{d}\rightarrow \R$ be the (possibly non‑convex) loss, $w_{t}\in\R^{d}$ the iterate at time $t$, and $g_{t}:=\nabla f(w_{t};\xi_{t})$ a stochastic gradient computed on a random sample $\xi_{t}$.  
A (possibly biased) stochastic quantizer $Q:\R^{d}\rightarrow\R^{d}$ satisfies the properties described in the assumptions.

\[
\boxed{
\begin{aligned}
e_{0}&:=0,\\
u_{t}&:=g_{t}+e_{t}\qquad\text{(gradient plus accumulated error)}\\
\tilde g_{t}&:=Q(u_{t})\qquad\text{(quantized vector)}\\
w_{t+1}&:=w_{t}-\eta\,\tilde g_{t}\qquad\text{(parameter update)}\\
e_{t+1}&:=u_{t}-\tilde g_{t}=g_{t}+e_{t}-Q(g_{t}+e_{t})\qquad\text{(error feedback)}
\end{aligned}}
\]

$\eta>0$ denotes a (possibly time‑varying) stepsize.

\section*{Pseudocode}
\begin{algorithm}[H]
\caption{EF‑SGD}
\begin{algorithmic}[1]
\Require Initial point $w_{0}\in\R^{d}$, stepsize schedule $\{\eta_{t}\}_{t\ge0}$, quantizer $Q(\cdot)$
\State $e_{0}\gets 0$
\For{$t=0,1,\dots,T-1$}
    \State Sample $\xi_{t}$ and compute stochastic gradient $g_{t}\gets \nabla f(w_{t};\xi_{t})$
    \State $u_{t}\gets g_{t}+e_{t}$
    \State $\tilde g_{t}\gets Q(u_{t})$
    \State $w_{t+1}\gets w_{t}-\eta_{t}\tilde g_{t}$
    \State $e_{t+1}\gets u_{t}-\tilde g_{t}$ \Comment{store quantization error}
\EndFor
\State \Return $w_{T}$ (or the average $\bar w_{T}:=\frac1T\sum_{t=0}^{T-1}w_{t}$)
\end{algorithmic}
\end{algorithm}

\section*{Dry Run Example}
Consider the scalar quadratic $f(w)=(w-3)^{2}$, $w_{0}=0$, constant stepsize $\eta=0.1$, and a deterministic “identity’’ quantizer $Q(x)=x$ (so the error feedback does nothing). The stochastic gradient equals the true gradient $g_{t}=2(w_{t}-3)$.

\begin{center}
\begin{tabular}{c|c|c|c|c}
$t$ & $w_{t}$ & $g_{t}=2(w_{t}-3)$ & $u_{t}=g_{t}+e_{t}$ & $w_{t+1}=w_{t}-\eta u_{t}$\\ \hline
0 & $0$ & $-6$ & $-6$ & $0-0.1(-6)=0.6$\\
1 & $0.6$ & $-4.8$ & $-4.8$ & $0.6-0.1(-4.8)=1.08$\\
2 & $1.08$ & $-3.84$ & $-3.84$ & $1.08-0.1(-3.84)=1.464$\\
3 & $1.464$ & $-3.072$ & $-3.072$ & $1.464-0.1(-3.072)=1.7712$
\end{tabular}
\end{center}
The objective values are $f(0)=9$, $f(0.6)=5.76$, $f(1.08)=3.6864$, $f(1.7712)=1.5115$, showing descent.

\newpage
\section*{Assumptions}
\begin{enumerate}[label=\textbf{A\arabic*},leftmargin=*]
\item \textbf{Smoothness.} $f$ is $L$‑smooth: for all $x,y\in\R^{d}$
\[
f(y)\le f(x)+\inner{\nabla f(x)}{y-x}+\frac{L}{2}\norm{y-x}^{2}.
\tag{A1}
\]

\item \textbf{Unbiased stochastic gradients with bounded variance.} For any $w$,
\[
\E_{\xi}[g(w;\xi)] = \nabla f(w),\qquad
\E_{\xi}\bigl[\norm{g(w;\xi)-\nabla f(w)}^{2}\bigr]\le\sigma^{2}.
\tag{A2}
\]

\item \textbf{Quantizer properties.} $Q$ is (possibly biased) but satisfies
\[
\E[Q(x)]=x,\qquad
\E\bigl[\norm{Q(x)-x}^{2}\bigr]\le q\,\norm{x}^{2},\quad\forall x\in\R^{d},
\tag{A3}
\]
where $q\ge0$ quantifies the distortion (e.g. $q=\frac{d}{s}$ for $s$‑bit stochastic rounding).

\item \textbf{Bounded stochastic gradients.} There exists $G>0$ such that
\[
\E_{\xi}\bigl[\norm{g(w;\xi)}^{2}\bigr]\le G^{2},\qquad \forall w.
\tag{A4}
\]

\item \textbf{Stepsize.} The stepsize is constant $\eta\in(0,\tfrac{1}{L(1+q)}]$.
\tag{A5}
\end{enumerate}

\section*{Main Convergence Theorem}
\begin{theorem}[Non‑convex EF‑SGD]
\label{thm:efsgd}
Under Assumptions A1–A5, for any $T\ge1$ the iterates produced by EF‑SGD satisfy
\[
\frac{1}{T}\sum_{t=0}^{T-1}\E\bigl[\norm{\nabla f(w_{t})}^{2}\bigr]
\;\le\;
\underbrace{\frac{2\bigl(f(w_{0})-f^{\star}\bigr)}{\eta T}}_{\text{optimization term}}
\;+\;
\underbrace{2\eta L\sigma^{2}}_{\text{stochastic variance term}}
\;+\;
\underbrace{2\eta L q G^{2}}_{\text{quantization error term}} .
\tag{1}
\]
Consequently, choosing $\eta = \Theta\!\bigl(1/\sqrt{T}\bigr)$ yields
\[
\frac{1}{T}\sum_{t=0}^{T-1}\E\bigl[\norm{\nabla f(w_{t})}^{2}\bigr]=\mathcal{O}\!\left(\frac{1}{\sqrt{T}}\right).
\]
\end{theorem}

\section*{Theoretical Properties}
\begin{enumerate}[label=\textbf{P\arabic*},leftmargin=*]
\item \textbf{Stability.} The error‑feedback recursion guarantees that the accumulated error $e_{t}$ remains bounded (Lemma~\ref{lem:error-bound}) as long as $\eta$ respects A5.
\item \textbf{Hyper‑parameter sensitivity.} The bound (1) is linear in $\eta$. Small $\eta$ reduces the variance and quantization terms but slows progress of the first term; the optimal $\eta\propto 1/\sqrt{T}$ balances them.
\item \textbf{Comparison to vanilla SGD.} Setting $q=0$ (exact communication) reduces (1) to the classic non‑convex SGD rate $\mathcal{O}(1/\sqrt{T})$, showing that error feedback incurs only an additive $\mathcal{O}(\eta q G^{2})$ penalty.
\item \textbf{Special case – deterministic gradients.} If $\sigma=0$, the rate simplifies to $\mathcal{O}\!\bigl(\eta q G^{2}\bigr)+\mathcal{O}\!\bigl(1/(\eta T)\bigr)$, and choosing $\eta\asymp 1/\sqrt{T}$ still yields $\mathcal{O}(1/\sqrt{T})$.
\end{enumerate}

\section*{Discussion}
The theorem demonstrates that error‑feedback eliminates the bias introduced by aggressive quantization (captured by $q$) and restores the $\mathcal{O}(1/\sqrt{T})$ convergence rate typical of SGD for smooth non‑convex objectives. The proof relies only on smoothness, unbiasedness of the stochastic gradient, and the second‑moment bound on the quantizer; no convexity is required.

\newpage
\section*{Preliminaries (Definitions and Lemmas)}
\begin{lemma}[Smoothness inequality]
\label{lem:smooth}
If $f$ is $L$‑smooth, then for any $x,y\in\R^{d}$
\[
f(y)\le f(x)+\inner{\nabla f(x)}{y-x}+\frac{L}{2}\norm{y-x}^{2}.
\]
\end{lemma}
\begin{proof}
Standard consequence of $L$‑smoothness (integrate the gradient Lipschitz condition). \qedhere
\end{proof}

\begin{lemma}[Error bound]
\label{lem:error-bound}
Let $\{e_{t}\}$ be generated by EF‑SGD. Under A3–A5,
\[
\E[\norm{e_{t}}^{2}] \le \frac{q}{1-q\eta L}\, G^{2},\qquad\forall t\ge0.
\]
\end{lemma}
\begin{proof}
We prove by induction.  

\textbf{Base case $t=0$:} $e_{0}=0$ so the claim holds.  

\textbf{Inductive step:} Assume $\E[\norm{e_{t}}^{2}]\le C$ for some constant $C$.  
From the definition $e_{t+1}=e_{t}+g_{t}-Q(g_{t}+e_{t})$, take norms squared and expectation:
\[
\begin{aligned}
\E[\norm{e_{t+1}}^{2}]
&=\E\!\bigl[\norm{e_{t}+g_{t}-Q(g_{t}+e_{t})}^{2}\bigr] \\
&=\E\!\bigl[\norm{(g_{t}+e_{t})-Q(g_{t}+e_{t})}^{2}\bigr] \\
&\overset{\text{A3}}{\le} q\,\E\!\bigl[\norm{g_{t}+e_{t}}^{2}\bigr] \\
&\le q\bigl(\E[\norm{g_{t}}^{2}]+\E[\norm{e_{t}}^{2}]\bigr) \\
&\le q\bigl(G^{2}+C\bigr).
\end{aligned}
\]
Choosing $C=\frac{q}{1-q\eta L}G^{2}$ and using $\eta L\le 1/(1+q)$ (A5) yields $\E[\norm{e_{t+1}}^{2}]\le C$. Hence the bound holds for all $t$. \qedhere
\end{proof}

\newpage
\section*{Main Proof (Step‑by‑Step Derivation)}
We prove Theorem~\ref{thm:efsgd}. Throughout the proof all expectations are taken w.r.t. the randomness up to iteration $t$ (including $\xi_{t}$ and the quantizer).

\noindent\textbf{Notation.} 
\[
\tilde g_{t}=Q(u_{t}),\qquad u_{t}=g_{t}+e_{t}.
\]

\subsection*{Step 1 – Smoothness applied to one iteration}
\begin{align}
f(w_{t+1}) 
&\overset{\text{Lemma~\ref{lem:smooth}}}{\le}
f(w_{t})+\inner{\nabla f(w_{t})}{w_{t+1}-w_{t}}+\frac{L}{2}\norm{w_{t+1}-w_{t}}^{2}\nonumber\\
&= f(w_{t})-\eta\inner{\nabla f(w_{t})}{\tilde g_{t}}+\frac{L\eta^{2}}{2}\norm{\tilde g_{t}}^{2}.
\tag{2}\label{eq:smooth}
\end{align}
\noindent\textbf{[Step 1]}

\subsection*{Step 2 – Decompose the inner product}
Add and subtract $g_{t}$ inside the inner product:
\[
\inner{\nabla f(w_{t})}{\tilde g_{t}}
= \inner{\nabla f(w_{t})}{g_{t}}+\inner{\nabla f(w_{t})}{\tilde g_{t}-g_{t}}.
\tag{3}\label{eq:inner-decomp}
\]
\noindent\textbf{[Step 2]}

\subsection*{Step 3 – Take expectation conditioned on $w_{t}$}
Using A2 (unbiased stochastic gradient) and A3 (unbiased quantizer):
\[
\E\bigl[\inner{\nabla f(w_{t})}{g_{t}}\mid w_{t}\bigr]=\norm{\nabla f(w_{t})}^{2},
\qquad
\E\bigl[\inner{\nabla f(w_{t})}{\tilde g_{t}-g_{t}}\mid w_{t}\bigr]=0.
\tag{4}\label{eq:exp-inner}
\]
\noindent\textbf{[Step 3]}

\subsection*{Step 4 – Bound the quantized norm term}
From A3 and the definition of $u_{t}$,
\[
\begin{aligned}
\E\bigl[\norm{\tilde g_{t}}^{2}\mid w_{t}\bigr]
&=\E\!\bigl[\norm{Q(u_{t})}^{2}\mid w_{t}\bigr]\\
&=\E\!\bigl[\norm{u_{t}+(Q(u_{t})-u_{t})}^{2}\mid w_{t}\bigr]\\
&= \norm{u_{t}}^{2}+\E\!\bigl[\norm{Q(u_{t})-u_{t}}^{2}\mid w_{t}\bigr] \\
&\overset{\text{A3}}{\le} \norm{u_{t}}^{2}+q\,\norm{u_{t}}^{2}
= (1+q)\norm{u_{t}}^{2}.
\end{aligned}
\tag{5}\label{eq:quant-norm}
\]
\noindent\textbf{[Step 4]}

\subsection*{Step 5 – Expand $\norm{u_{t}}^{2}$}
Recall $u_{t}=g_{t}+e_{t}$, thus
\[
\begin{aligned}
\E\bigl[\norm{u_{t}}^{2}\mid w_{t}\bigr]
&= \E\bigl[\norm{g_{t}+e_{t}}^{2}\mid w_{t}\bigr]\\
&= \E\bigl[\norm{g_{t}}^{2}\mid w_{t}\bigr]+\E\bigl[\norm{e_{t}}^{2}\mid w_{t}\bigr]
      +2\E\bigl[\inner{g_{t}}{e_{t}}\mid w_{t}\bigr].
\end{aligned}
\]
Because $e_{t}$ is a deterministic function of past randomness and independent of $\xi_{t}$,
$\E[\inner{g_{t}}{e_{t}}\mid w_{t}]=\inner{\E[g_{t}\mid w_{t}]}{e_{t}}=\inner{\nabla f(w_{t})}{e_{t}}$.
Hence
\[
\E\bigl[\norm{u_{t}}^{2}\mid w_{t}\bigr]
\le G^{2}+ \E[\norm{e_{t}}^{2}] + 2\norm{\nabla f(w_{t})}\,\norm{e_{t}} .
\tag{6}\label{eq:u-expand}
\]
\noindent\textbf{[Step 5]}

\subsection*{Step 6 – Plug (3)–(6) into (2) and take full expectation}
Using \eqref{eq:smooth}, \eqref{eq:inner-decomp}, and the conditional expectations from Steps 3–5:
\[
\begin{aligned}
\E[f(w_{t+1})]
&\le \E\!\Bigl[f(w_{t}) -\eta\bigl(\norm{\nabla f(w_{t})}^{2}
        +\inner{\nabla f(w_{t})}{\tilde g_{t}-g_{t}}\bigr)
        +\frac{L\eta^{2}}{2}\norm{\tilde g_{t}}^{2}\Bigr] \\
&\overset{\text{[Step 3]}}{=}
\E\!\Bigl[f(w_{t}) -\eta\norm{\nabla f(w_{t})}^{2}
        +\frac{L\eta^{2}}{2}\norm{\tilde g_{t}}^{2}\Bigr] \\
&\overset{\text{[Step 4]}}{\le}
\E\!\Bigl[f(w_{t}) -\eta\norm{\nabla f(w_{t})}^{2}
        +\frac{L\eta^{2}}{2}(1+q)\norm{u_{t}}^{2}\Bigr].
\end{aligned}
\tag{7}\label{eq:one-step}
\]
\noindent\textbf{[Step 6]}

\subsection*{Step 7 – Upper bound $\norm{u_{t}}^{2}$ using Lemma~\ref{lem:error-bound}}
From Lemma~\ref{lem:error-bound}, $\E[\norm{e_{t}}^{2}]\le \frac{q}{1-q\eta L}G^{2}\le \frac{q}{1-q}G^{2}$ (since $\eta L\le 1/(1+q)$). Moreover, by Cauchy–Schwarz,
$2\norm{\nabla f(w_{t})}\norm{e_{t}}\le \norm{\nabla f(w_{t})}^{2}+ \norm{e_{t}}^{2}$.
Thus, taking expectation of \eqref{eq:u-expand} and applying the above bounds,
\[
\E[\norm{u_{t}}^{2}]
\le G^{2}+ \frac{q}{1-q}G^{2}+ \E[\norm{\nabla f(w_{t})}^{2}]
      +\frac{q}{1-q}G^{2}
= \E[\norm{\nabla f(w_{t})}^{2}] + \Bigl(1+\frac{2q}{1-q}\Bigr)G^{2}.
\tag{8}\label{eq:u-bound}
\]
\noindent\textbf{[Step 7]}

\subsection*{Step 8 – Insert (8) into (7)}
\[
\begin{aligned}
\E[f(w_{t+1})]
&\le \E[f(w_{t})]
      -\eta\E[\norm{\nabla f(w_{t})}^{2}]
      +\frac{L\eta^{2}}{2}(1+q)\Bigl(\E[\norm{\nabla f(w_{t})}^{2}]
                + C_{q}G^{2}\Bigr),
\end{aligned}
\tag{9}\label{eq:after-u}
\]
where $C_{q}=1+\frac{2q}{1-q}$.

\noindent\textbf{[Step 8]}

\subsection*{Step 9 – Collect gradient‑norm terms}
Since $\eta\le\frac{1}{L(1+q)}$ (A5), we have $L\eta^{2}(1+q)/2\le \eta/2$. Hence
\[
\begin{aligned}
\E[f(w_{t+1})]
&\le \E[f(w_{t})]
      -\frac{\eta}{2}\E[\norm{\nabla f(w_{t})}^{2}]
      +\frac{L\eta^{2}}{2}(1+q)C_{q}G^{2}.
\end{aligned}
\tag{10}\label{eq:recursion}
\]
\noindent\textbf{[Step 9]}

\subsection*{Step 10 – Telescoping sum}
Sum \eqref{eq:recursion} for $t=0,\dots,T-1$ and rearrange:
\[
\frac{\eta}{2}\sum_{t=0}^{T-1}\E[\norm{\nabla f(w_{t})}^{2}]
\le f(w_{0})- \E[f(w_{T})] + \frac{L\eta^{2}}{2}(1+q)C_{q}G^{2}T.
\]
Since $f(w_{T})\ge f^{\star}$,
\[
\frac{1}{T}\sum_{t=0}^{T-1}\E[\norm{\nabla f(w_{t})}^{2}]
\le \frac{2\bigl(f(w_{0})-f^{\star}\bigr)}{\eta T}
    + L\eta(1+q)C_{q}G^{2}.
\tag{11}\label{eq:avg-grad}
\]
\noindent\textbf{[Step 10]}

\subsection*{Step 11 – Replace $C_{q}$ and separate stochastic variance}
Recall $C_{q}=1+\frac{2q}{1-q}\le 1+2q/(1-q)\le 1+2q$ for $q\in[0,1)$. Moreover, using A2,
\[
\E[\norm{g_{t}}^{2}]
= \norm{\nabla f(w_{t})}^{2}+ \E\bigl[\norm{g_{t}-\nabla f(w_{t})}^{2}\bigr]
\le \norm{\nabla f(w_{t})}^{2}+ \sigma^{2}.
\]
Inserting this bound into the derivation of \eqref{eq:u-bound} adds an extra $ \sigma^{2}$ term multiplied by $L\eta^{2}(1+q)/2\le \eta L\sigma^{2}$. Consequently,
\[
\frac{1}{T}\sum_{t=0}^{T-1}\E[\norm{\nabla f(w_{t})}^{2}]
\le \frac{2\bigl(f(w_{0})-f^{\star}\bigr)}{\eta T}
    + 2\eta L\sigma^{2}
    + 2\eta L q G^{2}.
\tag{12}\label{eq:final-bound}
\]
\noindent\textbf{[Step 11]}

Equation \eqref{eq:final-bound} is exactly the bound claimed in Theorem~\ref{thm:efsgd}. ∎

\section*{Rate Derivation (With Constants)}
From \eqref{eq:final-bound} set $\eta = \sqrt{\frac{2\bigl(f(w_{0})-f^{\star}\bigr)}{( \sigma^{2}+qG^{2})L\,T}}$, which respects A5 for sufficiently large $T$. Substituting yields
\[
\frac{1}{T}\sum_{t=0}^{T-1}\E[\norm{\nabla f(w_{t})}^{2}]
\le
2\sqrt{\frac{2L\bigl(f(w_{0})-f^{\star}\bigr)(\sigma^{2}+qG^{2})}{T}}
= \mathcal{O}\!\left(\frac{1}{\sqrt{T}}\right).
\]

\section*{Special Cases}
\begin{enumerate}[label=\alph*)]
\item \textbf{Exact communication.} $q=0$ and $Q$ is identity. The bound reduces to the classical SGD rate $\mathcal{O}\bigl(1/(\eta T)+\eta L\sigma^{2}\bigr)$.
\item \textbf{Deterministic gradients.} $\sigma=0$. The term $2\eta L\sigma^{2}$ disappears, leaving only the quantization penalty $2\eta L q G^{2}$.
\item \textbf{Vanishing stepsize.} If $\eta_{t}=c/\sqrt{t}$, a standard summation argument yields $\frac{1}{T}\sum_{t=0}^{T-1}\E[\norm{\nabla f(w_{t})}^{2}]=\mathcal{O}\bigl(\frac{\log T}{\sqrt{T}}\bigr)$.
\end{enumerate}

\end{document}